\documentclass{resume} % Use the custom resume.cls style

\usepackage[left=0.4 in,top=0.4in,right=0.4 in,bottom=0.4in]{geometry}
\usepackage{hanging}
\usepackage{enumitem}
\usepackage[hidelinks]{hyperref}
\usepackage{verbatim}

\name{Amber X. Chen} % Your name
\address{\href{mailto:amber.chen@psych.ucsb.edu}{amber.chen@psych.ucsb.edu} \\
\href{https://www.linkedin.com/in/amber-x-chen-b49057a9/}{LinkedIn} \\
\href{https://amberxchen.github.io/}{Homepage} \\
\href{https://github.com/amberxuqianchen}{GitHub} \\
\href{https://orcid.org/0000-0003-3359-9075}{ORCID}} \\

\begin{document}

%----------------------------------------------------------------------------------------
% OBJECTIVE
%----------------------------------------------------------------------------------------
\begin{rSection}{OBJECTIVE}
Mater student in Computer Science and PhD candidate in Social Psychology, seeking a \textbf{Data Scientist Intern} role for Summer 2026. 
Specialized in machine learning, NLP and applications of langguage and multi-modal models. 
Proven track record analysing large-scale data and building predictive models with statistical rigor.
\end{rSection}

%----------------------------------------------------------------------------------------
% EDUCATION
%----------------------------------------------------------------------------------------
\begin{rSection}{EDUCATION}
{\bf University of California, Santa Barbara (UCSB)} \hfill {\em 09/2022 -- (expected 05/2027)} \\
PhD Candidate in Social Psychology \hfill Santa Barbara, CA, USA 

{\bf University of California, Santa Barbara (UCSB)} \hfill {\em 09/2025 -- (expected 05/2027)} \\
Master's in Computer Science \hfill Santa Barbara, CA, USA 

{\bf The Chinese University of Hong Kong (CUHK)} \hfill {\em 11/2017} \\
Master's in Psychology \hfill Hong Kong SAR, China

{\bf Sun Yat-sen University (SYSU)} \hfill {\em 06/2016} \\
Bachelor's in Sociology \hfill Guangdong, China
\end{rSection}

%----------------------------------------------------------------------------------------
% SKILLS
%----------------------------------------------------------------------------------------
\begin{rSection}{SKILLS}
\textbf{Programming \& Data:} Python, SQL, R, MATLAB; Pandas, NumPy, SciPy, statsmodels \\
\textbf{Machine Learning \& NLP:} BERT, GPT, CLIP, VGG, PyTorch, TensorFlow, Scikit-learn, Hugging Face \\
\textbf{Statistical Methods:} A/B Testing, Causal Inference, Regression Analysis, Time Series, Bayesian Inference \\
\textbf{Data Infrastructure:} Google Cloud Platform, BigQuery, PySpark, SQL databases, APIs \\
\textbf{Analytics \& Visualization:} Statistical modeling, forecasting, Matplotlib, Seaborn, Tableau

\end{rSection}

%----------------------------------------------------------------------------------------
% PROJECTS
%----------------------------------------------------------------------------------------
\begin{rSection}{PROJECTS}

\noindent\textbf{Information Spread on Social Media}  — 
\href{https://osf.io/b7rjk/overview?view_only=569f284163104f37a1b5f5af43d9386b}{[Code]} \ \href{https://www.tandfonline.com/doi/full/10.1080/1369118X.2023.2246551}{[Paper]}
\hfill \textit{Current}
\begin{itemize}[leftmargin=*, noitemsep, topsep=2pt]
    \item \textbf{Analyzed 20M+ multilingual social media posts} on X and Weibo using advanced NLP pipelines (BERT, translation APIs) to track misinformation spread and user behavior patterns across cultural contexts
    \item \textbf{Built predictive models} to identify viral content patterns, achieving 85\% accuracy in predicting post engagement

\end{itemize}

\noindent\textbf{faceGPT: Multimodal AI for User Perception}\hfill \textit{Current}
\begin{itemize}[leftmargin=*, noitemsep, topsep=2pt]
    \item \textbf{Processed} 10 year+ of web news texts
    \item \textbf{Fine-tuned GPT/VGG/CLIP and evaluated} alignments between AI representations and behavioral data
    \item \textbf{Developed dual-route model} integrating language models to predict user impressions on visual content (images, logos, artwork), improving accuracy by +20\% over baseline models
\end{itemize}

\noindent\textbf{Time-Series Analysis of Cultural Sentiment} — 
\href{https://github.com/amberxuqianchen/effort-osf}{[Code]} \ 
\href{https://doi.org/10.1057/s41599-024-03603-3}{[Paper]} \ 
\href{https://www.nature.com/articles/s41599-024-03603-3.pdf}{[PDF]}
\hfill \textit{09/2023}
\begin{itemize}[leftmargin=*, noitemsep, topsep=2pt]
    \item \textbf{Scraped 70+ years of news articles}
    \item \textbf{Trained language models on 3.5B-token} using High-Performance Computing
    \item \textbf{Built time-series forecasting models and applied causal inference methods} on sentiment trends
    \item \textbf{Developed reproducible ML pipeline} with automated data processing, model training, and evaluation
\end{itemize}

\noindent\textbf{Deep Learning Health Prediction} — 
\href{https://doi.org/10.18632/aging.204264}{[Paper]}
\href{https://www.aging-us.com/news-room/Psychological-factors-substantially-contribute-to-biological-aging-evidence-from-the-aging-rate-in-Chinese-older-adults}{[Press Release]}
\hfill \textit{07/2022}  
\begin{itemize}[leftmargin=*, noitemsep, topsep=2pt]
    \item \textbf{Designed deep neural networks} for longitudinal health data analysis, achieving mean absolute error of 5.68 years on aging prediction
    \item \textbf{Implemented feature engineering pipeline} processing medical datasets with 2,000+ participants

\end{itemize}

\end{rSection}
%----------------------------------------------------------------------------------------

\begin{rSection}{WORK	EXPERIENCE}

\item \textbf{ Data Scientist \& Lab Manager} \hfill {\em Sept 2017 -- July 2022} \\
Motivation and Emotion Laboratory, CUHK
\begin{itemize}[leftmargin=*, noitemsep, topsep=2pt]
\item \textbf{Managed cross-cultural research projects} analyzing emotional patterns using statistical modeling, psychophysiological data, and large-scale surveys across multiple countries
\item \textbf{Supervised team of 20+ research assistants} in data collection, processing, and analysis workflows
\item \textbf{Present findings to stakeholders} and publish research outcomes
\end{itemize}

\item \textbf{Data Analyst \& Research Assistant } \hfill {\em May 2014 -- August 2016} \\
Center for Social Survey, SYSU
\begin{itemize}[leftmargin=*, noitemsep, topsep=2pt]
\item \textbf{Managed large-scale data collection} for the Chinese Education Panel Survey (CEPS), coordinating data from 1,000+ students across multiple schools
\item \textbf{Processed and analyzed survey data} using Stata to identify educational trends and policy implications
\end{itemize}


\item \textbf{Data Analyst Intern} \hfill {\em Feb 2014 -- May 2014} \\
MaxInsight, Customer Experience Analytics Platform
\begin{itemize}[leftmargin=*, noitemsep, topsep=2pt]
\item \textbf{Analyzed} user comments to generate actionable insights for competing brands for enterprise clients
\end{itemize}

\end{rSection}
%----------------------------------------------------------------------------------------
\begin{rSection}{AWARDS AND HONORS}
\item African and Asian Students in STEM Scholarship, UCSB \hfill {\em 2025}
\item GSA Conference Travel Grant, UCSB \hfill {\em 2023 - 2025}
\item Highlighted Poster Award, The 2019 Technology, Mind \& Society Conference, Washington, D.C., USA. \hfill {\em 2019}
\item Academic Excellence Scholarship (top 5\%), School of Sociology and Anthropology, SYSU \hfill {\em 2013 - 2015} 
\item Third Place Award at Survey Contest of Political Affairs, Guangdong Province \hfill {\em 2014} \\

\end{rSection}
%----------------------------------------------------------------------------------------
% SKILLS
%----------------------------------------------------------------------------------------
\begin{rSection}{LANGUAGES}

\item English (fluent), Mandarin Chinese (native), Cantonese (conversational)
\end{rSection}
\end{document}
